% !TEX program = pdflatex
\documentclass[10pt,twocolumn]{article}

% ---------- Packages ----------
\usepackage[utf8]{inputenc}
\usepackage[T1]{fontenc}
\usepackage{times}
\usepackage[margin=0.75in]{geometry}
\usepackage{amsmath,amssymb,amsthm}
\usepackage{graphicx}
\usepackage{booktabs}
\usepackage{caption}
\usepackage{subcaption}
\usepackage{hyperref}
\usepackage{xcolor}
\usepackage{enumitem}
\usepackage{listings}
\usepackage{multirow}
\usepackage{array}
\usepackage{balance}
\usepackage{microtype}
\usepackage{url}

% Help LaTeX break tight two-column lines without overflowing.
\tolerance=1000
\emergencystretch=3em
\hbadness=10000
\vbadness=10000
\sloppy
% Allow \url{} to break at any character if needed.
\makeatletter
\g@addto@macro{\UrlBreaks}{\UrlOrds}
\makeatother

\hypersetup{
  colorlinks=true,
  linkcolor=blue!60!black,
  citecolor=blue!60!black,
  urlcolor=blue!60!black
}

\lstset{
  basicstyle=\ttfamily\small,
  breaklines=true,
  frame=single,
  columns=fullflexible,
  keepspaces=true,
  xleftmargin=1em,
  framexleftmargin=0.5em
}

\newtheorem{theorem}{Theorem}
\newtheorem{proposition}[theorem]{Proposition}
\newtheorem{definition}{Definition}

% Compact sections
\usepackage{titlesec}
\titlespacing*{\section}{0pt}{1.5ex plus .5ex minus .2ex}{0.8ex plus .2ex}
\titlespacing*{\subsection}{0pt}{1.2ex plus .5ex minus .2ex}{0.6ex plus .2ex}
\titlespacing*{\subsubsection}{0pt}{0.8ex plus .5ex minus .2ex}{0.4ex plus .2ex}

\setlength{\textfloatsep}{8pt plus 2pt minus 2pt}
\setlength{\floatsep}{8pt plus 2pt minus 2pt}
\setlength{\intextsep}{6pt plus 2pt minus 2pt}
\setlength{\abovedisplayskip}{6pt plus 2pt minus 2pt}
\setlength{\belowdisplayskip}{6pt plus 2pt minus 2pt}

\title{\textbf{V-Lake: A Verifiable Lakehouse Architecture\\
for Multi-Party Governed Data}}

\author{
  Srinath Murali \and Rishav Kumar\\[4pt]
  \textit{Department of Computer Science}\\
  \textit{VIT University, Vellore, India}
}
\date{}

\begin{document}
\maketitle

% ============================================================
\begin{abstract}
Many regulated settings require several organisations to share
structured data without any one of them holding unilateral
authority over access, schema, or audit. Clinical trials, financial
trade reporting, cross-institution scientific collaborations, and
sustainability (ESG) disclosure are four examples. Existing data
lakehouse platforms address storage and query federation but
leave governance, integrity verification, and column-level privacy
to external tooling.

We present \textbf{V-Lake}, an architecture for verifiable
multi-party data governance over federated lakehouses. V-Lake is
\emph{not} a new cryptographic primitive. It is a specific
composition of four mechanisms that, to our knowledge, have not
previously been integrated at the lakehouse layer:
(i)~a domain-separated Merkle forest with on-chain root anchoring
and a per-query chained compliance attestation;
(ii)~a \emph{Weighted Quorum Governance} (WQG) layer---a voting
scheme atop the blockchain's QBFT ordering, not a consensus
algorithm---whose role weights are the unique minimal integer
solution to six explicit constraints and whose power distribution
we verify via Shapley--Shubik analysis;
(iii)~a DID-based consent chain providing subject-initiated
delegation and revocation (a practical subset of W3C SSI/VC, not a
full implementation);
and (iv)~a federated source registry spanning ten heterogeneous
systems. On top of these, a hybrid column-level encryption scheme
combines AES-SIV for near-unique identifiers with Fernet plus HMAC
blind indexes over \emph{deliberately coarsened} values for
clinical attributes.

We use a multi-site clinical trial as the driving case study
because it exercises all four pillars concurrently, but the
mechanisms are domain-agnostic; \S\ref{sec:applicability}
sketches how the same design applies to trade reporting, ESG
disclosure, and cross-institution research with only the
sensitive-column registry and policy pack changing. We validate
the design against a Hyperledger Besu (QBFT) network with a
830-line Solidity contract deployed and a Flask backend running
in \texttt{mode=on-chain}, exercised through a benchmark artefact
(\texttt{scripts/bench\_full.py}) that hits the real contract. Measured mean end-to-end query overhead over raw DuckDB
on five federated datasets is 10--21\% (approximately
0.3--1.0\,ms); Merkle build scales linearly from 0.06\,ms at
10 rows to 232\,ms at 50k rows; column-level encryption costs
0.03--0.08\,ms per value; a real on-chain ingestion anchor
round-trip (signed transaction plus block inclusion) takes a mean
of 23\,ms. Against a centralised IAM baseline (plain RBAC, no
chain, no Merkle), V-Lake adds roughly 1\,ms per query plus a
one-time per-ingestion blockchain cost of tens of milliseconds.
A test suite exercises seven modelled adversary classes; we do
not claim formal security proofs or regulatory certification. We
report the design, its rationale, its limitations, and the
tradeoffs we consciously accept.
\end{abstract}

% ============================================================
\section{Introduction}
\label{sec:intro}

Data lakehouses~\cite{armbrust2021lakehouse} consolidate the
schema-on-read flexibility of data lakes with the transactional
guarantees of data warehouses. In regulated multi-party settings,
however, a lakehouse must additionally provide \emph{data
integrity}, \emph{access governance}, \emph{subject
sovereignty}\footnote{Throughout, ``subject'' denotes the data
principal---a patient in healthcare, a beneficial owner in finance,
a disclosing entity in ESG---not a trial participant specifically.},
and \emph{privacy at rest}---properties that current architectures
leave to external tooling.

We use a multi-site clinical trial as the driving case study
because it exercises all four pillars concurrently: a pharmaceutical
sponsor, a hospital site, and an ethics board share patient
enrollment data (stored in S3), adverse-event reports (arriving via
Kafka), laboratory results (in PostgreSQL), vitals telemetry (in
MongoDB), and imaging documents (uploaded as PDFs). No single
organisation should be able to unilaterally modify governance
rules, grant access to sensitive columns, or alter data without
detection. Patients enrolled in the trial should retain control
over who sees their records, including the ability to delegate
access to a physician and revoke it at will, without requiring
approval from the governance layer.

The same four pillars appear with different names in other
regulated domains---multi-bank trade reporting under MiFID~II,
corporate ESG disclosure under the EU CSRD, and cross-institution
scientific data-use agreements among them. We return to this
point in \S\ref{sec:applicability}; the core architecture is not
healthcare-specific.

\subsection{What V-Lake is---and what it is not}

V-Lake is a principled composition of four mechanisms, not a new
cryptographic primitive. Each building block (Merkle trees,
weighted voting games, decentralised identifiers, searchable
encryption) is well-established. What is \emph{new}, to our
knowledge, is their specific integration at the lakehouse layer:

\begin{enumerate}[leftmargin=*,topsep=3pt,itemsep=2pt]
\item \textbf{Domain-Separated Merkle Forest (C1).}
  A typed hash tree with disjoint leaf/node domain tags and
  positional binding, extended with a cross-dataset forest root and
  a per-query hash-linked compliance attestation anchored on-chain.

\item \textbf{Weighted Quorum Governance (C2).}
  A governance voting scheme layered \emph{above} the blockchain's
  QBFT ordering (it does not replace consensus). Role weights are
  derived as the unique minimal integer solution to six stated
  constraints (Proposition~\ref{prop:unique}); the resulting power
  distribution is cross-checked via Shapley--Shubik analysis.

\item \textbf{DID-Based Consent Chain (C3).}
  An ECDSA-signed, hash-linked consent log using decentralised
  identifiers as subject handles. It provides subject-initiated
  delegation and revocation but is a \emph{subset} of the W3C
  SSI/VC model; we do not claim full VC-JSON interoperability
  (\S\ref{sec:ssi}).

\item \textbf{Federated Data Source Registry (C4).}
  Ten connectors, each labelled as \emph{materialising} the source
  into DuckDB or supporting predicate push-down
  (Table~\ref{tab:sources}).
\end{enumerate}

\noindent
On top of these, a hybrid column-level encryption scheme applies
deterministic encryption for near-unique identifiers, randomised
encryption for clinical values, and HMAC blind indexes over
deliberately coarsened buckets when filterability is required.

The specific technical contributions are:
(a)~a \emph{constraint-driven} derivation of role weights rather
than parameters chosen by convention, together with a full
Shapley--Shubik power-index check;
(b)~a cross-dataset Merkle forest whose root is tied, via a
hash-linked attestation chain, to the \emph{injected} SQL of each
query (not merely to ingested rows); and
(c)~a per-column encryption mode assignment driven by column
cardinality rather than applied uniformly.

We deliberately do not claim: formal security proofs, regulatory
certification, resistance to adversaries outside the stated threat
model, or full W3C VC interoperability. The evaluation is a
prototype datapoint against a single multi-site trial scenario,
not a broad performance claim.

The remainder of the paper proceeds as follows.
Section~\ref{sec:prior} surveys prior art.
Section~\ref{sec:model} defines the system and threat models and
provides a threat-mechanism matrix (Table~\ref{tab:threat-matrix}).
Section~\ref{sec:arch} presents the architecture.
Sections~\ref{sec:merkle}--\ref{sec:colcrypto} detail each
component. Section~\ref{sec:access} covers access-control
mechanics. Section~\ref{sec:eval} reports evaluation, methodology,
and threats to validity.
Section~\ref{sec:applicability} sketches how the architecture
applies beyond healthcare---to trade reporting, ESG disclosure,
and cross-institution research---and what would change per domain.
Section~\ref{sec:discussion} discusses what is genuinely new versus
inherited from prior work.
Sections~\ref{sec:limitations}--\ref{sec:conclusion} conclude.


% ============================================================
\section{Prior Art}
\label{sec:prior}

\paragraph{Data lakehouses.}
Delta Lake~\cite{armbrust2020delta}, Apache Iceberg~\cite{iceberg},
and Apache Hudi~\cite{hudi} provide ACID transactions and schema
evolution over object-store data but delegate governance and access
control to external catalogue services. V-Lake extends the lakehouse
model with built-in verifiable governance rather than competing with
these storage formats.

\paragraph{Blockchain-anchored data integrity.}
Anchoring Merkle roots on a blockchain is well-established in
certificate transparency~\cite{laurie2014ct} and supply-chain
provenance~\cite{toyoda2017supply}. MedRec~\cite{azaria2016medrec}
applies Ethereum smart contracts to medical-record permission
management. V-Lake builds on this insight but targets the lakehouse
layer rather than EMR systems, and adds a cross-dataset forest root
and per-query attestation chain absent in prior work.

\paragraph{Governance and weighted voting.}
Weighted voting has been studied extensively in cooperative game
theory~\cite{shapley1954method}. Systems such as Apache
Ranger~\cite{ranger} and Open Policy Agent~\cite{opa} enforce
role-based policies but do not expose formal power-index analysis
of their role hierarchies, and they execute unilaterally (a single
administrator can change a policy). V-Lake's WQG layer differs on
both points: weights are derived from explicit constraints, and
critical operations require a multi-party quorum recorded on-chain.

\paragraph{Self-sovereign identity.}
Decentralised identifiers (DIDs)~\cite{w3cdid} and verifiable
credentials~\cite{w3cvc} provide the foundation for
subject-controlled access. V-Lake uses a DID-style subject handle
and an ECDSA-signed consent chain but does not implement the full
Verifiable Credentials data model or DIDComm messaging; we are
explicit about this scoping in \S\ref{sec:ssi}.

\paragraph{Searchable encryption.}
CryptDB~\cite{popa2011cryptdb} introduced onion encryption with
deterministic and order-preserving layers for SQL. Subsequent work
showed that order-preserving encryption leaks enough information to
recover plaintexts on real medical data~\cite{naveed2015inference}.
V-Lake avoids OPE entirely, using deterministic encryption only for
near-unique identifiers and HMAC blind indexes with deliberate
coarsening for clinical values.


% ============================================================
\section{System and Threat Model}
\label{sec:model}

\subsection{Participants}

V-Lake defines five abstract roles with strictly ordered privilege
levels:
\textsc{Data~Steward}~$>$~\textsc{Data~Custodian}~$>$%
~\textsc{Analyst}~$>$~\textsc{Subject}~$>$~\textsc{None}.
The names reflect the clinical-trial case study; the contract
accepts role-label metadata so a deployment can surface
domain-appropriate labels (e.g., \textsc{Issuer}/\textsc{Auditor}/%
\textsc{Investor}/\textsc{Disclosing-Entity} for ESG,
\textsc{Desk}/\textsc{Compliance}/\textsc{Analyst}/\textsc{Customer}
for finance) while the weights and quorum rules remain unchanged.
Multiple organisations each contribute stewards; no single
organisation holds a majority.

\subsection{Trust Assumptions}

\begin{itemize}[leftmargin=*,topsep=3pt,itemsep=2pt]
\item The permissioned blockchain (Hyperledger Besu, QBFT) is
  trusted for ordering and availability but not for
  confidentiality.
\item At least one steward is honest. As shown in
  \S\ref{sec:wqg}, this suffices to block any critical governance
  change.
\item The application server is trusted for computation but not
  for long-term storage: integrity is verified via Merkle proofs
  against on-chain roots.
\item The query engine (DuckDB, in-memory) does not persist to
  disk, but we do not assume protection against OS-level memory
  attacks (swap, core dumps). Column-level encryption provides
  partial mitigation.
\end{itemize}

\subsection{Threat Model}

We consider seven adversary classes; Table~\ref{tab:threat-matrix}
maps each class to the V-Lake mechanism that is \emph{intended}
to address it, and to the empirical evidence presented in
\S\ref{sec:eval}. We stress that this is a design-and-test mapping,
not a formal security proof.

\begin{table*}[t]
\centering
\footnotesize
\caption{Threat classes, V-Lake mechanisms intended to address them,
and the evidence presented. ``Test-validated'' means a concrete
adversarial scenario is implemented in our test suite and is
observed to be rejected; it does not imply the mechanism is secure
against all adversaries in the class.}
\label{tab:threat-matrix}
\setlength{\tabcolsep}{4pt}
\begin{tabular}{@{}c p{3.6cm} p{4.0cm} p{2.8cm} p{3.6cm}@{}}
\toprule
\textbf{\#} & \textbf{Adversary class} & \textbf{Primary mechanism}
  & \textbf{Secondary} & \textbf{Evidence (\S\ref{sec:eval})} \\
\midrule
T1 & Data tampering (storage/transit)
   & Domain-separated Merkle forest (C1)
   & On-chain anchor
   & Test-validated: root mismatch \\
T2 & Governance subversion (rogue coalition)
   & Weighted Quorum Governance (C2)
   & \texttt{all\_stew} unanimity
   & Test-validated: 2/3 critical op denied \\
T3 & Consent forgery (altered records)
   & SHA-256 hash-linked consent chain (C3)
   & ECDSA signatures
   & Test-validated: chain verify rejects \\
T4 & Unauthorised column access
   & Non-Truman predicate injection (\S\ref{sec:access})
   & Grant allow-list
   & Test-validated: HTTP 403 returned \\
T5 & SQL injection in row filter
   & Deny-list token validator
   & Parameter escaping
   & Test-validated; \emph{see caveat} \\
T6 & Stolen backup / rogue DBA (at rest)
   & Column-level encryption (\S\ref{sec:colcrypto})
   & Envelope key hierarchy
   & Ciphertext in DuckDB verified \\
T7 & Expired or revoked grant reuse
   & Temporal check + grant cache invalidation
   & Redis-backed shared cache
   & Test-validated: rejection observed \\
\bottomrule
\end{tabular}
\end{table*}

\noindent
\emph{Caveat on T5.} The current row-filter validator is a
deny-list, which is known to be weaker than a SQL-parser-based
allow-list and may miss novel attack vectors. We treat this as a
known limitation (\S\ref{sec:limitations}), not a solved problem.

\emph{Caveat on T6.} The column-encryption scheme has documented
frequency-leakage tradeoffs on blind-indexed columns
(\S\ref{sec:colcrypto}). ``Encrypted at rest'' is not the same as
``information-theoretically secure.''

Adversaries \emph{outside} these seven classes---for example,
side-channel attacks on the host OS, a dishonest majority of
stewards, or a compromise of the Besu network's validator set---are
not addressed by this work.


% ============================================================
\section{Architecture Overview}
\label{sec:arch}

Figure~\ref{fig:arch} illustrates V-Lake's pipeline; subsequent
sections detail each component.

\subsection{Pipeline}

Data enters V-Lake through one of ten federated connectors
(\S\ref{sec:federation}). Each source is either materialised into
DuckDB's in-memory columnar engine or accessed via push-down
(Table~\ref{tab:sources}). The pipeline then proceeds as follows:

\begin{enumerate}[leftmargin=*,topsep=3pt,itemsep=2pt]
\item \textbf{Column encryption.}
  PHI columns are identified by a default registry (overridable per
  dataset) and encrypted in place: deterministic encryption for
  identifiers, randomised encryption for clinical values, with
  blind-index side columns for coarsened equality filters
  (\S\ref{sec:colcrypto}).

\item \textbf{Merkle commitment.}
  A domain-separated Merkle tree is built over the
  \emph{post-encryption} rows, producing a root $\rho$ that commits
  to the ciphertext state (\S\ref{sec:merkle}).

\item \textbf{On-chain anchoring.}
  $\rho$ is recorded on a Hyperledger Besu smart contract (QBFT,
  chain~ID~1337, 2\,s block period). A forest root $\Phi$
  aggregates all dataset roots into a single verifiable hash.

\item \textbf{Governance.}
  Role assignments, policy attachments, and access grants are
  mediated by WQG (\S\ref{sec:wqg}). Every governance action is
  recorded on-chain with a quorum certificate.

\item \textbf{Query processing.}
  Incoming queries pass through predicate injection (column and row
  restriction per the caller's grant), PHI predicate rewriting
  (translating plaintext literals to ciphertext or blind-index
  values), compliance checking (HIPAA, GDPR, DPDPA rules), and
  Merkle-proof generation. Results are decrypted for authorised
  callers only.
\end{enumerate}

\begin{figure}[t]
\centering
\scriptsize
\begin{verbatim}
 +-- Federated Sources (10 types) --+
 | MinIO  Postgres  Kafka  Mongo .. |
 +---------------+------------------+
                 |
           DuckDB (in-memory)
                 |
        +--------+---------+
        | Column Encrypt.  |
        | det/rand/bidx    |
        +--------+---------+
                 |
        +--------+---------+
        | Merkle Forest C1 |
        | domain-separated |
        +--------+---------+
                 |
        +--------+---------+
        | Besu Contract    |
        | anchor + audit   |
        +--------+---------+
                 |
  +--------------+---------------+
  |              |               |
 WQG (C2)   SSI Chain(C3)  Compliance
 governance  consent log   HIPAA/GDPR
  |              |               |
  +--------------+---------------+
                 |
          Query Pipeline
       predicate injection
       PHI rewrite + decrypt
       Merkle proof bundle
                 |
              Caller
\end{verbatim}
\caption{V-Lake architecture. Data flows top to bottom; governance
and consent are cross-cutting concerns checked at query time.}
\label{fig:arch}
\end{figure}

\subsection{Technology Stack}

The prototype is implemented as a Flask (Python) backend with
DuckDB as the query engine, Hyperledger Besu as the permissioned
blockchain (QBFT, 4 validators in reference config; a single-node
dev instance in the benchmark), and MinIO, PostgreSQL, Kafka, and
MongoDB as real federated sources. The smart contract
(\textasciitilde830\,LoC Solidity) mirrors the backend's governance
state on-chain. The system is deployed via Docker Compose
(8~containers). All measurements in \S\ref{sec:eval} are taken
with the backend in \texttt{mode=on-chain}; measurements are
invalidated if that flag is absent.


% ============================================================
\section{C1: Domain-Separated Merkle Forest}
\label{sec:merkle}

\subsection{Construction}

We use SHA-256 with two domain-separated hash functions to prevent
second-preimage confusion between leaves and internal nodes:
%
\begin{align}
H_{\text{leaf}}(r, i) &= \textsc{Sha256}\!\bigl(
  \texttt{"vlake.leaf:"} \,\|\, i \,\|\, \texttt{":"} \,\|\,
  n_c \,\|\, \texttt{":"} \,\|\, \hat{r}\bigr) \\[2pt]
H_{\text{node}}(l, r, \ell) &= \textsc{Sha256}\!\bigl(
  \texttt{"vlake.node:"} \,\|\, \ell \,\|\, \texttt{":"} \,\|\,
  l \,\|\, \texttt{":"} \,\|\, r\bigr)
\end{align}
%
\noindent
Here $i$ is the row's zero-based positional index, $n_c$ is the
column count, $\hat{r}$ is the canonical serialisation of the row,
and $\ell$ is the tree level.

\subsection{Design Properties}

Three properties follow from this construction:
\emph{domain separation} ensures leaf and node hashes cannot
collide under SHA-256 collision resistance;
\emph{position binding} (the index $i$ appears in the leaf hash)
prevents row-reorder attacks;
\emph{odd-leaf promotion} (unpaired nodes are promoted unchanged,
not duplicated) prevents a length-extension variant where
$\text{root}(A,B,C) = \text{root}(A,B,C,C)$.
The empty sentinel root is
$H_\emptyset = \textsc{Sha256}(\texttt{"vlake.empty\_tree"})$.

\subsection{Forest Root and On-Chain Anchoring}

For each dataset $j$ with Merkle root $\rho_j$, a forest leaf is
$f_j = \textsc{Sha256}(\texttt{"vlake.forest:"} \| j \| \texttt{":"} \| \rho_j)$,
and $\Phi$ is the Merkle root over the $\{f_j\}$. After every
ingestion, V-Lake calls
$\textsc{RecordIngestion}(\text{id}, \rho, |R|, |L|, \text{depth})$
on the Besu contract; the contract maintains a per-dataset
\texttt{merkleHistory[id]} and the global $\Phi$.

Each query additionally produces a chained attestation
$a_t = \textsc{Sha256}(H(q_t) \| h_t \| \rho_t \| a_{t-1} \| \tau_t)$,
where $H(q_t)$ is the \emph{injected}-query hash (i.e., after
predicate injection), $h_t$ is the result hash, $\rho_t$ the
dataset root, $a_{t-1}$ the previous attestation, and $\tau_t$ the
timestamp. This chain, anchored on-chain, is the property that
distinguishes V-Lake's integrity layer from a plain Merkle log: the
SQL that produced each answer---not only the stored rows---is part
of the tamper-evident record.

\subsection{Document Ciphertext Anchoring}

The per-dataset Merkle root $\rho$ commits to the rows of a
dataset's \emph{table} in the query engine, including extracted
text and metadata for document uploads. It does \emph{not} by
itself commit to the raw ciphertext object stored in the
object-store backend (e.g., a MinIO bucket, S3, or IPFS). A rogue
storage admin who swaps an object with different ciphertext of the
same length leaves the Merkle root unchanged.

We close this gap without a separate on-chain record and without
changing the integrity-proof protocol: when a document $d$ is
ingested, we compute
$\kappa_d = \textsc{Sha256}(\text{ciphertext}(d))$
and store $\kappa_d$ as an additional column
(\texttt{ciphertext\_hash}) in the document row. Because $\kappa_d$
is part of the row, it is part of the leaf hash
$H_{\text{leaf}}(r, i)$ (Eq.~1), and therefore part of the
anchored root $\rho$. At download time the server refetches the
ciphertext, recomputes $\kappa'_d$, and refuses to decrypt when
$\kappa'_d \neq \kappa_d$, logging the mismatch as a
\texttt{CIPHERTEXT\_TAMPER\_DETECTED} audit event. The expected
$\kappa_d$ ultimately traces to the on-chain root, so an attacker
who tampers with both the storage object and the in-memory
metadata is detected on any subsequent Merkle verification.

The property this gives us is that object-store durability and
object-store integrity become \emph{separable} concerns: the
storage tier can be any addressable backend (MinIO, S3,
erasure-coded cluster, IPFS pins) without affecting the integrity
guarantee, because integrity is rooted in the chain, not the
store.


% ============================================================
\section{C2: Weighted Quorum Governance}
\label{sec:wqg}

\subsection{Terminology: governance, not consensus}

V-Lake's weighted voting layer is a \emph{governance} scheme, not a
\emph{consensus} algorithm. Blockchain consensus (Besu's QBFT)
orders transactions; WQG determines \emph{which} governance
transactions are ratified. Earlier drafts of this work used
``Weighted Quorum Consensus'' (WQC); we have retitled it to avoid
confusion. Internal identifiers in our implementation retain the
historical \texttt{WQC} prefix for backwards compatibility with
the deployed ABI.

\subsection{Weight Derivation}

V-Lake assigns integer weights: Steward~$s{=}3$, Custodian~$c{=}2$,
Analyst~$a{=}1$. We state six requirements and prove that
$(3,2,1)$ is the unique minimal integer solution.

\begin{enumerate}[label=\textbf{R\arabic*},leftmargin=2em,
  topsep=3pt,itemsep=2pt]
\item \textbf{Steward supremacy}: two stewards alone meet the
  standard quorum, i.e., $2s \geq q$.
\item \textbf{No unilateral steward}: one steward alone cannot
  pass, i.e., $s < q$.
\item \textbf{Custodian relevance}: a custodian is pivotal in at
  least one winning coalition.
\item \textbf{Custodian insufficiency}: non-stewards alone cannot
  pass, i.e., $c + a < q$.
\item \textbf{Strict ordering}: $0 < a < c < s$.
\item \textbf{Integer weights}: $s, c, a \in \mathbb{Z}^+$ (for
  gas-efficient on-chain representation).
\end{enumerate}

\noindent
Fix the reference committee used throughout: three stewards, one
custodian, one analyst, total $W = 3s + c + a$ and standard quorum
$q = \lceil W/2 \rceil$.

\begin{proposition}\label{prop:unique}
Under \textbf{R1}--\textbf{R6} applied to the reference committee,
$(s, c, a) = (3, 2, 1)$ is the unique minimal integer solution,
where minimality is with respect to the lexicographic order
$(s, c, a)$.
\end{proposition}

\begin{proof}
By \textbf{R5}--\textbf{R6}, $a \geq 1$, $c \geq a+1 \geq 2$, and
$s \geq c+1 \geq 3$. We show $(3,2,1)$ satisfies all six
constraints and that no lexicographically smaller integer triple
does.

\smallskip
\emph{Feasibility of $(3,2,1)$.}
$W = 9 + 2 + 1 = 12$; $q = 6$.
\textbf{R1}: $2s = 6 \geq 6$.~\checkmark\
\textbf{R2}: $s = 3 < 6$.~\checkmark\
\textbf{R4}: $c + a = 3 < 6$.~\checkmark\
\textbf{R3}: the coalition $\{S_1, S_2, C_1\}$ has weight
$3+3+2 = 8 \geq 6$ (winning), but removing $C_1$ leaves
$\{S_1, S_2\}$ with weight $6 \geq 6$ (still winning), so $C_1$
is \emph{not} pivotal there. Consider instead $\{S_1, C_1, A_1\}$:
weight $3+2+1 = 6 \geq 6$ (winning); $\{S_1, A_1\}$ has weight
$4 < 6$ (losing), so $C_1$ is pivotal.~\checkmark

\smallskip
\emph{Minimality.}
By \textbf{R5}--\textbf{R6} the lex-smallest candidate is
$(s,c,a) = (3,2,1)$ itself; any strictly smaller triple must reduce
$s$, $c$, or $a$, which violates \textbf{R5}--\textbf{R6}.

\smallskip
\emph{No feasible triple of equal $s$ with smaller $(c,a)$.}
Fix $s = 3$. Then $W = 9 + c + a$, $q = \lceil W/2 \rceil$. For
$(c,a) = (2,1)$ we showed \textbf{R1}--\textbf{R4} hold. Consider
\emph{any} $(c', a')$ with $c' < 2$ or $a' < 1$: \textbf{R5}
forces $c' \geq 2$ and $a' \geq 1$, so no strictly smaller
$(c,a)$ exists subject to \textbf{R5}.

\smallskip
Combining, $(3,2,1)$ satisfies \textbf{R1}--\textbf{R6} and no
lex-smaller integer triple does. Uniqueness of the minimal
solution follows.
\end{proof}

\begin{proposition}[Safety: no double ratification]
\label{prop:safety}
For any quorum threshold $q \geq \lceil W/2 \rceil$, a proposal
cannot be simultaneously approved ($\textsc{Yes}\,\geq\,q$) and
rejected ($\textsc{No}\,\geq\,q$).
\end{proposition}

\begin{proof}
$\textsc{Yes} + \textsc{No} \leq W$. If $\textsc{Yes} \geq q$, then
$\textsc{No} \leq W - q \leq W - \lceil W/2 \rceil < q$ whenever
$W \geq 2$, so $\textsc{No} < q$.
\end{proof}

\begin{proposition}[Single-honest-steward sufficiency]
\label{prop:honest}
For critical operations that require \texttt{all\_stew}, a
single honest steward suffices to block ratification regardless of
the behaviour of all other participants.
\end{proposition}

\begin{proof}
Ratification requires \texttt{stewardsMet} $\equiv$ every steward
voted \textsc{Yes}. A single \textsc{No} from an honest steward
falsifies this predicate.
\end{proof}

\subsection{Shapley--Shubik Verification}

To cross-check that the induced power distribution matches the
intended hierarchy, we compute the Shapley--Shubik power
index~\cite{shapley1954method} for the reference game
$[7;\; 3, 3, 3, 2, 2, 1]$ (three stewards, two custodians, one
analyst) by full enumeration of all $6! = 720$ permutations
(Table~\ref{tab:shapley}). The ordering
$\varphi(S) > \varphi(C) > \varphi(A)$ matches the weight ordering
and the intended liability hierarchy.

\begin{table}[t]
\centering
\small
\caption{Shapley--Shubik power index for $[7; 3,3,3,2,2,1]$.}
\label{tab:shapley}
\begin{tabular}{@{}lccc@{}}
\toprule
\textbf{Role} & \textbf{Weight} & \textbf{Pivots\,/\,720}
  & $\boldsymbol{\varphi}$ \\
\midrule
Steward   & 3 & 156 & 0.217 \\
Custodian & 2 & 108 & 0.150 \\
Analyst   & 1 &  36 & 0.050 \\
\bottomrule
\end{tabular}
\end{table}

\subsection{Quorum Types}

V-Lake defines three quorum types (Table~\ref{tab:quorum}):
\emph{standard} (50\% weighted threshold), \emph{critical} (67\%
plus \texttt{all\_stew}), and \emph{emergency} (any single steward,
for HIPAA~\S164.312 incident response).

\begin{table}[t]
\centering
\small
\caption{Quorum configuration per governance operation.}
\label{tab:quorum}
\begin{tabular}{@{}lcccc@{}}
\toprule
\textbf{Operation} & \textbf{Type} & \textbf{Thr.}
  & \textbf{All\,$S$} & \textbf{$C$\,maj.} \\
\midrule
\textsc{AssignCustodian}    & std   & 50\% & -- & -- \\
\textsc{OnboardAnalyst}     & std   & 50\% & \checkmark & \checkmark \\
\textsc{AccessGrant}        & std   & 50\% & \checkmark & \checkmark \\
\textsc{RevokeCustodian}    & std   & 50\% & -- & -- \\
\textsc{RevokeAnalyst}      & emerg & 1\,$S$ & -- & -- \\
\textsc{AttachPolicy}       & crit  & 67\% & \checkmark & -- \\
\textsc{ToggleConfidential} & crit  & 67\% & \checkmark & \checkmark \\
\bottomrule
\end{tabular}
\end{table}


% ============================================================
\section{C3: DID-Based Consent Chain}
\label{sec:ssi}

\subsection{Scope and caveat}

V-Lake's consent layer uses decentralised identifiers as subject
handles and an ECDSA-signed, hash-linked log as the record of
consent. This provides the \emph{property} often associated with
SSI---that a data subject can unilaterally delegate or revoke
access without intermediation---but it is not a full W3C
SSI/VC deployment. In particular:

\begin{itemize}[leftmargin=*,topsep=3pt,itemsep=2pt]
\item DIDs are derived locally as
  $\text{DID} = \texttt{"did:vlake:"} \,\|\, \textsc{Sha256}%
  (\text{address} \| \tau)_{[0{:}16]}$;
  we do not register a DID method with the W3C.
\item Consent records are JSON documents with a fixed schema; they
  are not VC-JSON or JWT-VC credentials.
\item There is no DIDComm messaging layer; delegation is issued
  through the V-Lake API, not peer-to-peer.
\end{itemize}

\noindent
Interoperability with other SSI stacks is future work
(\S\ref{sec:limitations}). The consent-chain design gives subjects
the specific \emph{agency} properties listed below, within the
V-Lake ecosystem.

\subsection{Consent Chain Structure}

Each consent record $c_i$ carries:
\texttt{prev\_hash} $= c_{i-1}.\texttt{hash}$
(or \texttt{"genesis"} for $c_0$);
\texttt{hash} $= \textsc{Sha256}(\text{JSON}(c_i))$ with
deterministic key ordering;
\texttt{signature} $=$ ECDSA signature over the concatenation of
action, dataset, filter, timestamp, and \texttt{prev\_hash}.
Chain verification asserts
$c_i.\texttt{prev\_hash} = c_{i-1}.\texttt{hash}$ at every step; any
modification breaks all subsequent links.

\subsection{Subject-Initiated Operations}

Three operations require no governance-layer approval:
\emph{Link} (a steward binds a subject to a dataset with a SQL row
filter); \emph{Delegate} (a subject grants time-bounded access to
a delegate, default 86\,400\,s); and \emph{Revoke} (a subject
unilaterally deactivates a delegate's grant and invalidates the
query-engine cache immediately). This design gives patients agency
over their records without requiring steward intervention for
day-to-day delegation.


% ============================================================
\section{C4: Federated Data Source Registry}
\label{sec:federation}

V-Lake supports ten source types (Table~\ref{tab:sources}).
\textbf{Mode} indicates whether V-Lake materialises the source into
DuckDB at ingestion time or issues pushed-down predicates at query
time.

\begin{table}[t]
\centering
\footnotesize
\caption{Federated source types, connection method, and execution
mode.}
\label{tab:sources}
\setlength{\tabcolsep}{3pt}
\begin{tabular}{@{}p{2.3cm} p{3.3cm} p{1.5cm}@{}}
\toprule
\textbf{Source} & \textbf{Method} & \textbf{Mode} \\
\midrule
Local File (CSV/JSON) & DuckDB \texttt{read\_csv\_auto} & materialise \\
S3 / MinIO            & DuckDB \texttt{httpfs}          & materialise \\
PostgreSQL            & DuckDB \texttt{postgres\_scan}  & push-down* \\
MySQL / MariaDB       & DuckDB \texttt{mysql\_scan}     & push-down* \\
Trino / Starburst     & Trino JDBC bridge               & push-down \\
Delta Lake            & DuckDB \texttt{delta\_scan}     & materialise \\
Apache Iceberg        & DuckDB \texttt{iceberg\_scan}   & materialise \\
HDFS / Hadoop         & \texttt{httpfs} + WebHDFS       & materialise \\
Kafka                 & \texttt{kafka-python} consumer  & materialise \\
MongoDB               & \texttt{pymongo} export         & materialise \\
\midrule
Documents (PDF/DOCX/IMG)
  & Text extract + envelope enc.
  & materialise \\
\bottomrule
\end{tabular}
\end{table}

\smallskip
\noindent
\emph{* Push-down caveat.} DuckDB's scan extensions push simple
filter and projection predicates but materialise the remainder;
column-level encryption therefore runs on the materialised
partition, not at the source. Sources labelled ``push-down'' are
scalable beyond memory for filtered queries but degrade to
``materialise'' for unfiltered scans.

After materialisation (or push-down evaluation), every source goes
through the same pipeline: column encryption, Merkle tree
construction, and on-chain anchoring. Cross-source SQL joins are
supported because all sources resolve to tables in the same DuckDB
instance.

Documents follow an additional path: text is extracted (pypdf for
PDF, XML parsing for DOCX), the raw binary is encrypted with the
dataset's envelope key (Fernet: AES-128-CBC + HMAC-SHA256), and
the ciphertext is persisted to a MinIO bucket. Plaintext bytes
never survive beyond the upload handler.

The key design decision is \emph{materialisation by default, with
push-down where the source supports it}: a uniform point at which
encryption, Merkle hashing, and access control can be applied
consistently. The tradeoff is memory-bounded dataset size for
materialised sources; we discuss this in \S\ref{sec:limitations}.


% ============================================================
\section{Column-Level PHI Encryption}
\label{sec:colcrypto}

\subsection{Motivation}

HIPAA's Breach Safe Harbor provision applies only to encrypted
PHI. V-Lake's document layer already encrypts file uploads at rest,
but the structured-data path (DuckDB tables) previously stored
rows in plaintext. We address this gap with a hybrid scheme that
preserves SQL filterability where possible. We note explicitly
that encryption at rest is a \emph{necessary but not sufficient}
condition for Safe Harbor eligibility; regulatory determination
requires independent audit (\S\ref{sec:limitations}).

\subsection{Key Hierarchy}

A master KEK wraps per-dataset DEKs. Three subkeys are derived per
DEK via HKDF-like expansion:
%
\begin{align}
k_{\text{siv}}    &= \textsc{Hkdf}\bigl(\text{DEK},\;
  \texttt{"vlake.colcrypt.siv.v1"},\; 64\bigr) \\
k_{\text{fernet}} &= \textsc{Hkdf}\bigl(\text{DEK},\;
  \texttt{"vlake.colcrypt.fernet.v1"},\; 32\bigr) \\
k_{\text{bidx}}   &= \textsc{Hkdf}\bigl(\text{DEK},\;
  \texttt{"vlake.colcrypt.bidx.v1"},\; 32\bigr)
\end{align}

\subsection{Three Encryption Modes}

\textbf{Deterministic (det).}
AES-SIV~\cite{rogaway2006siv} with a 512-bit key. Identical
plaintexts always produce identical ciphertexts, enabling equality
filters and joins. Applied to \emph{high-cardinality identifiers}
(MRN, SSN, email, patient name) where frequency leakage is minimal.

\textbf{Randomised (rand).}
Fernet (AES-128-CBC + HMAC-SHA256) with a fresh IV per call.
Applied to \emph{clinical values} (diagnosis, DOB, ZIP, notes).

\textbf{Blind index (bidx).}
Truncated HMAC-SHA256 over a \emph{coarsened} value, stored in a
side column (\texttt{\_bidx\_\textit{col}}). Enables equality
filters with bounded, documented frequency leakage. Three
coarseners are provided: \texttt{year} (DOB $\to$ year),
\texttt{zip3} (first 3 ZIP chars), \texttt{icd3} (first 3
ICD-10 chars, e.g., E11.9 $\to$ E11).

A column that is \texttt{rand}-encrypted \emph{and} has a blind
index gets random ciphertext on the stored value while the side
column enables coarse equality filters.
Table~\ref{tab:colmodes} summarises the assignment.

\begin{table}[t]
\centering
\footnotesize
\caption{Column encryption mode assignment by column type.}
\label{tab:colmodes}
\setlength{\tabcolsep}{3pt}
\begin{tabular}{@{}p{2.1cm} c c p{2.5cm}@{}}
\toprule
\textbf{Column Type} & \textbf{Mode} & \textbf{Blind Idx}
  & \textbf{Rationale} \\
\midrule
MRN, SSN, email, name & det & -- & Near-unique; flat freq. \\
DOB, birth date       & rand & year & Quasi-identifier \\
ZIP, postal code      & rand & zip3 & Quasi-identifier \\
Diagnosis, ICD code   & rand & icd3 & Disease-family bucket \\
Notes, free text      & rand & -- & No filtering needed \\
Non-PHI columns       & --   & -- & Plaintext \\
\bottomrule
\end{tabular}
\end{table}

\subsection{Query Rewriting}

Equality predicates on PHI columns are rewritten before reaching
DuckDB. \texttt{det} columns: the literal is replaced with its
deterministic ciphertext. \texttt{rand}+\texttt{bidx} columns: the
predicate targets the blind-index side column with the HMAC of the
coarsened value. \texttt{rand}-only columns: the predicate becomes
\texttt{FALSE} (we prefer an explicit zero-row result to a silent
match of every row). Range queries, \texttt{LIKE}, and function
calls on encrypted columns are not rewritten.

\subsection{Frequency-Leakage Budget}

We are explicit about what this scheme leaks.
Deterministic encryption on near-unique columns produces a flat
frequency distribution. Blind indexes produce a bucketed
distribution; an attacker can correlate bucket frequencies with
published prevalence data (the classic \cite{naveed2015inference}
failure mode), which is why the coarsening is chosen per column
(family-level for ICD) rather than left at the raw value.
\texttt{rand}-only columns leak no frequency information.
Blind-index keys are per-dataset, so compromising one dataset does
not break frequency analysis across the lake. This is a
\emph{documented tradeoff}, not a solved problem.


% ============================================================
\section{Access Control and Compliance}
\label{sec:access}

\subsection{Non-Truman Predicate Injection}

When a grant restricts a caller to a subset of columns or rows:

\begin{itemize}[leftmargin=*,topsep=3pt,itemsep=2pt]
\item \textbf{Column restriction}: \texttt{SELECT~*} is rewritten
  to an explicit column list. Queries naming unauthorised columns
  return HTTP~403 with the disallowed column list---the
  ``Non-Truman'' model~\cite{rizvi2004nontruman}, which rejects
  rather than silently filters.
\item \textbf{Row restriction}: the grant's SQL predicate is
  prepended as
  \texttt{WHERE~(\textit{filter})~AND~(\textit{original})}.
\end{itemize}

\subsection{SQL Injection Prevention (known-weak)}

User-supplied row filters are validated against a deny-list of
dangerous tokens. This is a \emph{known weaker} defence than a
SQL-parser-based allow-list and may miss novel attack forms. A
production deployment should integrate a parser; we flag this
explicitly in Table~\ref{tab:threat-matrix} (T5) and
\S\ref{sec:limitations}.

\subsection{Compliance Engine}

Policy rules are grouped into \emph{policy packs} that define the
set of sensitive columns and the rules to apply. The current
prototype ships healthcare-oriented packs (HIPAA, GDPR, India
DPDPA) but the pack format is a plain declarative registry;
finance (SOX, MiFID~II, PCI-DSS), disclosure (EU CSRD), and
payments (PCI-DSS) packs are straightforward additions that do not
change the engine. The engine checks four properties per query:
column exposure (informational), temporal validity, audit logging,
and minimum-necessary (which blocks \texttt{SELECT~*} on datasets
with sensitive columns). Each policy emits a chained attestation
whose hash feeds the next policy's input; the overall chain is
anchored on-chain. The word ``compliance'' in this section refers
to rule checks matching the pack's patterns, not to regulatory
certification.

\subsection{Grant Lifecycle}

Grants are created through WQG proposals; they carry an access
level (\textsc{ViewOnly}, \textsc{ViewDownload}, \textsc{FullAccess}),
optional column restriction, optional row filter, and expiration.
A TTL-based grant cache (30\,s, Redis-backed in multi-instance
deployments) avoids repeated policy lookups; revocation invalidates
the cache immediately.


% ============================================================
\section{Evaluation}
\label{sec:eval}

\subsection{Scope: Architecture vs.\ Prototype}
\label{sec:arch-vs-prototype}

Before presenting numbers we separate two layers that are easily
conflated in this kind of work and that deserve independent
scrutiny.

\emph{The architecture} is the set of mechanisms described in
\S\S\ref{sec:merkle}--\ref{sec:access}: domain-separated Merkle
forest with forest root and per-query attestation chain, weighted
quorum governance, DID-style consent chain, federated source
registry, hybrid column encryption, ciphertext-hash anchoring.
These are stated generally; they do not depend on any specific
deployment topology and do not bake in a specific object store,
key-management system, or blockchain validator set.

\emph{The prototype} is the specific instantiation we built to
measure the architecture: Flask backend + DuckDB query engine +
single-node Hyperledger Besu (QBFT, dev-mode) + single-node MinIO
+ a disk-resident master key. It makes several simplifying
assumptions that a production deployment would not make:

\begin{enumerate}[leftmargin=*,topsep=3pt,itemsep=2pt]
\item The Besu network is a single validator. The architecture
  requires a 4+ validator QBFT cluster for real liveness against
  validator failures. The prototype is sufficient to measure
  round-trip latency with a real contract but is not a
  highly-available ordering layer.
\item The object store is a single-node MinIO. The architecture
  is store-agnostic; the prototype is not fault-tolerant against
  volume loss.
\item The master key-encryption key lives in a file on disk. The
  architecture requires an HSM- or KMS-backed KEK; the prototype
  uses a local file so the demo works without external secret
  management.
\item The query engine is in-memory DuckDB. The architecture
  commits to an SQL engine that honours predicate injection and
  blind-index rewrites; the prototype uses DuckDB because its
  extension model makes the federated registry tractable.
\item Data volumes are small (8--50 rows per table at demo time)
  because the prototype exists to demonstrate mechanism behaviour,
  not peak throughput.
\end{enumerate}

Every reported number in \S\ref{sec:eval} describes the prototype.
Claims that extend to the architecture in general are flagged
explicitly.


\subsection{Methodology}

All measurements are taken against the real stack: backend running
with \texttt{mode=on-chain}, the 830-line Solidity contract
deployed to a Hyperledger Besu (QBFT, dev-mode) node, and the same
five federated datasets seeded by the demo walkthrough. The
benchmark driver is \texttt{scripts/bench\_full.py}; it is
reproducible and is the artefact behind every number in this
section. The driver aborts if the backend is not in on-chain mode,
preventing accidental numbers from a simulated or cache-only run.

For each operation we report $n$, mean, standard deviation, 95\%
confidence interval (normal approximation), and the p50/p95/p99
percentiles. Query microbenchmarks use $n=200$ after 5 warmup
iterations; Merkle benchmarks use $n{=}5$ per size; on-chain
round-trips use $n{=}10$ for anchor ingestion and $n{=}5$ for the
full propose-and-vote cycle owing to the much higher per-sample
cost. The measurement host is a Windows 11 laptop with the Besu,
backend, and datastore containers co-located via Docker Compose
over an internal bridge network; RPC calls therefore include only
localhost and intra-container latency. Real-network latency will
be strictly higher.

\subsection{Datasets and Representativeness}

Five datasets from the demo walkthrough are loaded from four
distinct backends (MinIO, PostgreSQL, MongoDB, local documents)
with row counts of 8, 0, 15, 30, and 1. These sizes are
representative of a \emph{per-site} clinical-trial extract at the
prototype stage; they are not a scale benchmark, and we report
Merkle scaling separately to address that axis. We do not claim
generalisation to production-scale data warehouses.

\subsection{End-to-End Query Overhead}

Table~\ref{tab:overhead} reports mean, p95, and relative overhead
of the full V-Lake pipeline over raw DuckDB for each dataset. The
V-Lake path includes predicate injection, compliance checking, PHI
predicate rewriting, Merkle-proof generation, and attestation
hashing; the baseline path is DuckDB alone.

\begin{table}[t]
\centering
\footnotesize
\caption{End-to-end query overhead: V-Lake (full pipeline) vs raw
DuckDB, per dataset. $n{=}200$ iterations after warmup.
Measurements from \texttt{scripts/bench\_full.py} on the running
on-chain stack.}
\label{tab:overhead}
\setlength{\tabcolsep}{4pt}
\begin{tabular}{@{}lcccc@{}}
\toprule
\textbf{Dataset} & \multicolumn{2}{c}{\textbf{Baseline (ms)}}
                 & \multicolumn{2}{c}{\textbf{V-Lake (ms)}} \\
\cmidrule(lr){2-3}\cmidrule(lr){4-5}
                 & avg & p95 & avg & p95 \\
\midrule
trial\_enrollment & 4.04 & 6.33 & 4.88 & 7.57 \\
lab\_results      & 3.03 & 4.36 & 3.33 & 4.49 \\
vitals\_stream    & 4.92 & 8.11 & 5.60 & 7.63 \\
imaging\_reports  & 3.87 & 6.81 & 4.44 & 6.96 \\
\midrule
\textbf{Mean}     & 3.97 & 6.40 & 4.56 & 6.66 \\
\multicolumn{5}{@{}l}{\textbf{Overhead:} +0.59\,ms avg,
  +0.26\,ms p95 ($\approx$ 10--21\%)} \\
\bottomrule
\end{tabular}
\end{table}

\subsection{Comparison to a Centralised IAM Baseline}

A natural question is what V-Lake \emph{costs} relative to a
conventional centralised stack. We report a best-case centralised
baseline: a Python dict-lookup RBAC check plus a raw DuckDB query
on a synthetic 1,000-row table, with no chain, no Merkle tree, no
compliance attestation, and an append-only in-memory audit list
(Table~\ref{tab:centralized}). A single-admin grant mutation is
essentially a dict write, measured at~$\sim0$\,ms.

\begin{table}[t]
\centering
\small
\caption{Centralised IAM baseline vs V-Lake increments. The
centralised baseline uses a synthetic table and an in-memory ACL,
so it is an \emph{upper bound} on what a conventional stack
achieves.}
\label{tab:centralized}
\begin{tabular}{@{}lcc@{}}
\toprule
\textbf{Operation} & \textbf{Centralised} & \textbf{V-Lake} \\
\midrule
Query path (mean)                      & 1.94\,ms  & 4.56\,ms      \\
Query path (p95)                       & 3.01\,ms  & 6.66\,ms      \\
Admin grant (single writer)            & $\sim0$\,ms & n/a         \\
WQG proposal + 2 on-chain votes (mean) & n/a       & 18.84\,ms     \\
Ingestion Merkle anchor (on-chain)     & n/a       & 23.37\,ms     \\
Tamper detection                       & none      & root mismatch \\
Multi-party separation of duties       & none      & WQG           \\
Subject-initiated revocation           & none      & C3            \\
Cross-instance audit integrity         & none      & on-chain      \\
\bottomrule
\end{tabular}
\end{table}

\noindent
\textbf{Takeaway.} Per query, V-Lake adds approximately 1--3\,ms
over a centralised RBAC path---a constant-factor cost dominated by
Merkle-proof generation and compliance hashing. The large costs
(tens of milliseconds) are \emph{per-ingestion} and
\emph{per-governance-action}, not per-query. In exchange, V-Lake
provides four properties that the centralised baseline has no way
to offer: integrity verification, separation of duties, subject
sovereignty, and cross-instance audit integrity. Whether that cost
is worth paying is a deployment decision; our claim is only that
the incremental cost is small and the extra properties are
measurable.

\subsection{Merkle Build Scaling}

Merkle construction is $O(n)$; Table~\ref{tab:merkle-scale} shows
the measured numbers match this shape. The 50k-row case builds in
232\,ms mean.

\begin{table}[t]
\centering
\small
\caption{Merkle forest build time vs row count ($n{=}5$ per size).}
\label{tab:merkle-scale}
\begin{tabular}{@{}rcc@{}}
\toprule
\textbf{Rows} & \textbf{Avg (ms)} & \textbf{P95 (ms)} \\
\midrule
10       & 0.06   & 0.07   \\
100      & 0.34   & 0.37   \\
1\,000   & 3.56   & 3.83   \\
5\,000   & 19.15  & 22.82  \\
10\,000  & 36.01  & 38.44  \\
50\,000  & 231.87 & 279.69 \\
\bottomrule
\end{tabular}
\end{table}

\subsection{Column Encryption Kernel}

The column-encryption primitives are fast enough that their cost is
dominated by Python overhead; they are not a bottleneck in the
query path (Table~\ref{tab:crypto-bench}).

\begin{table}[t]
\centering
\small
\caption{Column encryption kernel microbenchmark.}
\label{tab:crypto-bench}
\begin{tabular}{@{}lcc@{}}
\toprule
\textbf{Operation} & \textbf{Mean (ms)} & \textbf{P95 (ms)} \\
\midrule
Deterministic (AES-SIV) encrypt+decrypt & 0.03 & 0.05 \\
Randomised (Fernet) encrypt+decrypt     & 0.08 & 0.12 \\
Blind-index HMAC (icd3 coarsener)       & 0.00 & 0.01 \\
\bottomrule
\end{tabular}
\end{table}

\subsection{On-Chain Round-Trip}

Every anchor and every governance action is a real signed
transaction submitted to Besu; the benchmark waits for the
transaction receipt before recording the sample
(Table~\ref{tab:onchain}). Anchor ingestion (one TX) has a mean
round-trip of 23.37\,ms; a full proposal-plus-two-vote cycle
(three TXs) has a mean round-trip of 18.84\,ms. The two numbers
are close because Besu in dev-mode produces blocks on demand when
the transaction pool is non-empty rather than waiting strictly for
the 2\,s period; the QBFT configuration would dominate on a
production network with validators separated by a real network.

\begin{table}[t]
\centering
\small
\caption{Real on-chain round-trip (Besu QBFT, dev-mode, localhost).
Each sample includes signing, submission, and receipt wait.}
\label{tab:onchain}
\begin{tabular}{@{}lccc@{}}
\toprule
\textbf{Operation} & \textbf{$n$} & \textbf{Mean (ms)}
  & \textbf{P95 (ms)} \\
\midrule
Merkle-root anchor (1 TX)            & 10 & 23.37 & 32.00 \\
WQG proposal + 2 steward votes (3 TX) & 5 & 18.84 & 31.90 \\
\bottomrule
\end{tabular}
\end{table}

These are not simulated values. They are taken with the backend
calling the contract deployed at
\texttt{0xb9A2\ldots C0cc} on the running Besu node; removing the
Besu container causes the backend health endpoint to switch out of
on-chain mode and the benchmark to abort.

\subsection{Adversary Scenarios (Scoped)}

We implement and run seven concrete adversary scenarios
(Table~\ref{tab:attacks}). Each scenario is a deterministic unit
test: the benchmark calls the endpoint, observes whether the
adversary succeeds, and reports the outcome. Every scenario is
rejected by the tested defence in this configuration. We are
careful to say ``test-validated'' rather than ``provably secure'':
these tests exhaust neither the stated threat class nor the
adversary space outside the threat model. The row-filter validator
in particular is a deny-list and is flagged as a known-weak
defence (T5).

\begin{table}[t]
\centering
\footnotesize
\caption{Adversary scenarios and the defence that rejected them in
our test suite. ``Test-validated'' is not a security proof.}
\label{tab:attacks}
\setlength{\tabcolsep}{3pt}
\begin{tabular}{@{}p{3.3cm} c p{2.8cm}@{}}
\toprule
\textbf{Scenario} & \textbf{Rejected} & \textbf{Defence} \\
\midrule
SQL injection in row filter
  & \checkmark & Deny-list (weak; see \S\ref{sec:limitations}) \\
Unauthorised column access
  & \checkmark & Non-Truman rejection \\
Role escalation (analyst votes)
  & \checkmark & Role check at vote time \\
Expired grant reuse
  & \checkmark & Temporal check in query engine \\
Data tampering
  & \checkmark & Merkle root mismatch \\
Consent chain forgery
  & \checkmark & Hash-chain verification \\
2/3 stewards on critical op
  & \checkmark & \texttt{all\_stew} unanimity \\
\bottomrule
\end{tabular}
\end{table}

\subsection{Threats to Validity}

\emph{External validity.} The evaluation is a single clinical-trial
scenario with small per-dataset row counts. We present Merkle
scaling, crypto microbench, and on-chain round-trip separately to
give axis-specific evidence, but the end-to-end latency numbers
should be read as \emph{prototype datapoints}, not production
performance claims.

\emph{Measurement noise.} On datasets of 8--30 rows the raw DuckDB
query itself takes 3--5\,ms, so the V-Lake overhead (0.3--1.0\,ms
on average) is close to the measurement noise floor. We report
confidence intervals and percentile spreads so readers can assess
this directly; larger-row studies are a natural follow-up.

\emph{Internal validity.} Benchmarks run co-located with the
backend on localhost; real-network latency, contention from other
workloads, and Besu validator-to-validator propagation on a
production network will increase the numbers, particularly for the
on-chain round-trip. We report the dev-mode Besu figure honestly
but note that a 2\,s QBFT block period on a real network would
dominate those measurements.


% ============================================================
\section{Applicability Beyond Healthcare}
\label{sec:applicability}

The mechanisms in \S\S\ref{sec:merkle}--\ref{sec:colcrypto} do not
depend on healthcare semantics. The four pillars---multi-party
governance without unilateral authority, tamper-evident integrity,
subject-initiated consent, and column-level encryption with
bounded filterability---recur in several regulated domains. We
sketch three below and list the concrete changes each would
require, all of which are configuration changes rather than
architectural changes.

\paragraph{Financial trade reporting (MiFID~II, Dodd--Frank, SOX).}
Multiple banks, clearinghouses, and regulators share order and
execution data for which no single party should hold unilateral
audit authority. WQG maps to the existing committee structure
(trading desks, risk office, compliance, external regulator); the
\texttt{all\_stew} property supports the regulator-as-gating-signer
pattern. Column encryption modes map directly: account IDs are
near-unique (\texttt{det}), counterparty codes are low-cardinality
(\texttt{rand}+\texttt{bidx}), notional amounts need
\emph{bucketed} blind indexes (new coarsener: log-scale amount
buckets). SOX~\S404 and MiFID~II~Art.~25 audit requirements are
served by the per-query attestation chain. The main architectural
gap is throughput: trade capture is write-heavy, so batch anchoring
(one Merkle root per $N$ trades) is required to avoid making each
trade a blockchain transaction.

\paragraph{ESG / sustainability disclosure (EU CSRD, SEC climate).}
Issuers, auditors, investors, and regulators jointly rely on
disclosure data that is increasingly scrutinised for greenwashing.
V-Lake's tamper-evident integrity and per-query attestation address
the core integrity concern; WQG matches the
issuer/auditor/regulator triangle (stewards = issuer + auditor +
regulator; custodian = ESG data manager; analyst = investor
analyst). Only the column registry changes: emissions scope
codes, facility IDs, and supplier identifiers replace MRN/ICD
fields, and the \texttt{icd3} family coarsener generalises to
scope/category bucketing.

\paragraph{Cross-institution scientific collaboration.}
Multi-site consortia (e.g., genomics networks, particle-physics
collaborations, climate-model intercomparisons) manage
data-use-agreement-gated datasets that resemble clinical trials in
governance shape but usually have looser column-level secrecy and
much larger data volumes. WQG and the consent chain apply
unchanged; column encryption is often unnecessary except for
personally identifying fields. The per-query attestation chain has
an additional benefit here: it makes query-level
\emph{reproducibility} auditable, which is a recognised pain point
in open science.

\paragraph{Changes the model would need to target these domains.}
The architecture is stable; what changes is (i)~the \emph{policy
pack} and sensitive-column registry, (ii)~the set of blind-index
coarseners (add numeric bucketing and date-quarter for finance;
scope/category for ESG), (iii)~optional role aliases in the
on-chain metadata so the UI displays domain-appropriate labels
while the contract still operates on the
Steward/Cus\-todian/Analyst/Sub\-ject abstraction, and (iv)~batch
anchoring of Merkle roots for write-heavy ingestion. The weight
derivation in Proposition~\ref{prop:unique}, the Shapley--Shubik
verification, the Merkle construction, and the encryption kernel
are unchanged.

Two things do \emph{not} generalise for free.
First, \textbf{right-to-erasure} under CCPA and CSRD conflicts
with append-only on-chain anchoring; a crypto-shredding protocol
(per-subject DEK destruction) is needed to make deletion
cryptographically effective without rewriting history. This is
flagged in \S\ref{sec:limitations}.
Second, \textbf{sub-second-latency} use cases (IoT telemetry,
automotive event streams) stress the QBFT block period and would
need either off-chain batching with periodic anchoring or a
different ordering layer.


% ============================================================
\section{Discussion: What Is Genuinely New}
\label{sec:discussion}

It is worth being explicit about what in this work is new versus
standard composition of known primitives, and what is not. We
split the answer into three categories.

\paragraph{Inherited, unchanged.}
SHA-256 Merkle trees, ECDSA signatures, AES-SIV, HMAC-SHA256,
QBFT consensus, the Shapley--Shubik power index, the Non-Truman
access-control model, and the observation that order-preserving
encryption leaks too much on medical data. None of these are
contributions.

\paragraph{Adapted, with small modifications.}
The domain-separated Merkle hashing style is folklore from
transparency logs; our contribution here is consistent application
(including position binding and odd-leaf promotion) and the
aggregation of per-dataset roots into a forest root with its own
domain tag. The Fernet-plus-blind-index pattern exists in various
SDKs; our contribution is the \emph{cardinality-driven} mode
assignment and explicit coarsening budget.

\paragraph{New, to our knowledge, at the lakehouse layer.}
(i)~A constraint-driven derivation of weighted-voting weights
(Proposition~\ref{prop:unique}) rather than parameters chosen by
convention, paired with a full power-index check---most
RBAC/policy stacks use convention and do not publish either.
(ii)~A per-query chained attestation that binds the
\emph{post-injection} SQL to the Merkle root and anchors the
chain on-chain; this makes query-level tamper detection possible,
not just row-level.
(iii)~The specific composition of a cross-dataset forest root,
WQG governance, a DID-style consent chain, and hybrid column
encryption behind a single query pipeline.

What is explicitly \emph{not} new and we do not claim otherwise:
no new cryptographic primitive; no formal security proof; no
regulatory certification; no full W3C SSI/VC implementation; and
no generalisation of the evaluation beyond the prototype
configuration described.


% ============================================================
\section{Limitations and Future Work}
\label{sec:limitations}

\begin{enumerate}[leftmargin=*,topsep=3pt,itemsep=2pt]
\item \textbf{In-memory materialisation.}
  DuckDB runs in-memory, limiting materialised dataset size to RAM.
  Push-down sources (Postgres, MySQL, Trino) scale beyond this but
  degrade for unfiltered scans. Production deployments would
  benefit from spill-to-disk with encrypted swap or a persistent
  encrypted engine.

\item \textbf{Blind-index frequency leakage.}
  Coarsening reduces leaked information but does not eliminate it.
  Blind indexes on low-cardinality columns still reveal bucketed
  frequency distributions. Documented tradeoff, not a solved
  problem.

\item \textbf{Predicate rewriter scope.}
  The current rewriter supports only simple equality predicates on
  PHI columns. Range queries, \texttt{LIKE}, and functions on
  encrypted columns are not rewritten. Range queries would require
  order-revealing encryption, which we deliberately
  avoid~\cite{naveed2015inference}.

\item \textbf{Row-filter validation is deny-list based.}
  The deny-list is known weaker than a parser-based allow-list. A
  production deployment should integrate a SQL parser. This is
  flagged in Table~\ref{tab:threat-matrix} (T5).

\item \textbf{DID/SSI interoperability.}
  The consent chain is not a full W3C VC/DIDComm implementation
  (\S\ref{sec:ssi}). Interop with external SSI stacks is future
  work.

\item \textbf{No formal security proof.}
  Security claims are supported by test-validated scenarios
  (\S\ref{sec:eval}) rather than formal proofs. Symbolic or
  game-based analysis is future work.

\item \textbf{No regulatory certification.}
  Phrases such as ``HIPAA Safe Harbor'' and ``GDPR Art.~32'' refer
  to design alignment, not to third-party certification. Real
  deployment requires independent audit.

\item \textbf{Benchmark scope.}
  Numbers in \S\ref{sec:eval} reflect a single host, dev-mode
  Besu, and small datasets. Scaling to larger committees, larger
  policy sets, and real-network validators is future work.

\item \textbf{Right-to-erasure vs append-only anchors.}
  On-chain Merkle roots are immutable, but GDPR, CCPA, India
  DPDPA, and EU CSRD each grant data-subject deletion rights.
  The prototype supports logical deletion and key rotation; a
  robust \emph{crypto-shredding} protocol (per-subject DEK so
  destroying the key erases data cryptographically while the
  anchor remains) is planned future work and would be required
  before deployment in any domain where subject-deletion is
  enforceable.

\item \textbf{Write-heavy throughput.}
  Trade capture, IoT, and similar write-heavy workloads cannot
  anchor each record on-chain within a 2\,s QBFT block period.
  The intended extension is batch anchoring: buffer $N$ ingestions
  and anchor one aggregated Merkle root. The security argument
  carries over (the aggregate root commits to all leaves), but
  the latency between ingestion and anchor becomes a design
  parameter.
\end{enumerate}

\subsection{Staged Storage-Integrity Hardening}
\label{sec:stages}

The prototype assumptions in \S\ref{sec:arch-vs-prototype} imply a
production rollout path. Each stage below is orthogonal: a
deployment can adopt them in the order that matches its risk
profile. Stage A is already implemented in the prototype; Stages
B--E are design commitments whose implementation is out of scope
for this paper.

\begin{itemize}[leftmargin=*,topsep=3pt,itemsep=3pt]

\item \textbf{Stage A. Ciphertext content-addressing (done).}
  The architecture is already store-agnostic because every
  document row carries $\kappa_d = \textsc{Sha256}(\text{ciphertext}(d))$
  and $\kappa_d$ is part of the Merkle leaf (\S\ref{sec:merkle}).
  At download, the store is refetched, $\kappa'_d$ is recomputed,
  and any mismatch with the on-chain-anchored $\kappa_d$ refuses
  the decrypt and emits a \texttt{CIPHERTEXT\_TAMPER\_DETECTED}
  audit event. This closes the \emph{silent object-swap} threat
  without changing the integrity-proof protocol.

\item \textbf{Stage B. Replicated object store.}
  The prototype uses single-node MinIO. A production deployment
  replaces this with a MinIO erasure-coded cluster (4+ nodes, RS
  parity) or equivalent S3-class service. Because Stage A makes
  integrity independent of the store, the replication layer needs
  only to guarantee \emph{availability}---integrity verification
  remains a property of the on-chain state.

\item \textbf{Stage C. Decentralised object store (IPFS).}
  Multi-organisation deployments where no single party should
  hold the binaries benefit from a content-addressed, replicated
  store such as IPFS with a per-organisation pin set or a
  Filecoin-backed pinning service. The architectural interface
  does not change: $\kappa_d$ is now both the on-chain-anchored
  hash and the IPFS CID, collapsing two concepts into one. IPFS
  does not substitute for encryption---ciphertext is still what is
  published---and it introduces DHT-resolution latency that makes
  it a better fit for long-term archival than interactive queries.
  A hybrid deployment can keep MinIO for the hot path and IPFS
  for cross-organisation replication and cold archive.

\item \textbf{Stage D. HSM- or KMS-backed key management.}
  The prototype stores the master KEK as a local file
  (\texttt{backend/data/.master.key}) because the demo must work
  without external secret infrastructure. Production requires an
  HSM (FIPS 140-3 Level 3) or a managed KMS (HashiCorp Vault,
  AWS KMS, GCP KMS) wrapping the KEK. Per-dataset DEKs remain
  wrapped ciphertext in the existing \texttt{deks.json}; the only
  change is that the wrap/unwrap call is delegated to the HSM/KMS.
  This stage is load-bearing for any HIPAA Safe Harbor or
  FIPS-aligned deployment.

\item \textbf{Stage E. Multi-validator Besu.}
  The prototype runs Besu in single-node dev-mode, which is
  sufficient to measure real transaction round-trip latency but
  is not a highly available ordering layer. Production QBFT
  requires 4 validators to tolerate one Byzantine fault (or $3f+1$
  for tolerance of $f$). This is a deployment-topology change
  rather than an architectural change; the smart contract is
  unchanged.

\end{itemize}

None of these stages alters the security argument, the weight
derivation in Proposition~\ref{prop:unique}, or the frequency-leakage
budget of the column-encryption scheme. They change what the
prototype's numbers generalise to.


% ============================================================
\section{Conclusion}
\label{sec:conclusion}

V-Lake demonstrates that verifiable governance, tamper-evident
integrity, subject-initiated delegation, and privacy-preserving
column encryption can be composed into a lakehouse architecture
with small per-query overhead (approximately 1\,ms over a
centralised IAM baseline on our prototype) and tens of milliseconds
per ingestion anchor on a real blockchain. Although we use
multi-site clinical trials as the driving case study, the mechanisms
are domain-agnostic: \S\ref{sec:applicability} shows how the same
architecture serves financial trade reporting, ESG disclosure, and
cross-institution scientific collaboration, with policy-pack and
coarsener changes rather than architectural changes. The contribution is not
a new primitive but a specific composition: weights derived from
explicit constraints and verified via Shapley--Shubik, a Merkle
forest whose root is tied to per-query attestations, a DID-style
consent chain scoped to subject-initiated delegation, and a
column-encryption scheme with a documented frequency-leakage
budget.

We have been deliberate about what the paper does and does not
claim, and about the split between architecture and prototype
(\S\ref{sec:arch-vs-prototype}). The architecture is stated
generally; the prototype makes explicit simplifying assumptions
(single-validator chain, single-node object store, disk-resident
KEK) that \S\ref{sec:stages} organises into a staged hardening
plan, of which Stage A---ciphertext content-addressing with
verify-on-download---is already implemented. The paper does not
claim a new cryptographic primitive, a formal security proof,
regulatory certification, or generalisation beyond the prototype. It does claim that the composition is coherent, the
weight derivation is minimal and unique under the stated
constraints, and the measured overhead is small enough to be a
practical alternative to centralised IAM when the additional
properties (multi-party separation of duties, tamper-evident
integrity, subject sovereignty, cross-instance audit) are
required.

\paragraph{Reproducibility.} All numbers in this paper are produced
by \texttt{scripts/bench\_full.py} against the running stack with
the backend in \texttt{mode=on-chain}. The script aborts if the
stack is not live; removing the Besu container breaks it. This is
deliberate: we do not want the paper's numbers to be mistakable for
simulation output.


% ============================================================
\balance
\begin{thebibliography}{20}

\bibitem{armbrust2021lakehouse}
M.~Armbrust, A.~Ghodsi, R.~Xin, and M.~Zaharia,
``Lakehouse: A New Generation of Open Platforms that Unify Data
Warehousing and Advanced Analytics,''
in \emph{Proc.\ CIDR}, 2021.

\bibitem{armbrust2020delta}
M.~Armbrust \textit{et al.},
``Delta Lake: High-Performance ACID Table Storage over Cloud Object
Stores,''
\emph{Proc.\ VLDB Endow.}, vol.~13, no.~12, pp.~3411--3424, 2020.

\bibitem{iceberg}
R.~Blue, O.~Meara, and D.~McIntosh,
``Apache Iceberg: An Open Table Format for Huge Analytic Datasets,''
2020. \url{https://iceberg.apache.org}

\bibitem{hudi}
V.~Jain \textit{et al.},
``Apache Hudi: The Streaming Data Lake Platform,'' 2020.
\url{https://hudi.apache.org}

\bibitem{laurie2014ct}
B.~Laurie, A.~Langley, and E.~K\"asper,
``Certificate Transparency,'' RFC~6962, IETF, 2013.

\bibitem{toyoda2017supply}
K.~Toyoda, P.~T.~Mathiopoulos, I.~Sasase, and T.~Ohtsuki,
``A Novel Blockchain-Based Product Ownership Management System
(POMS) for Anti-Counterfeits in the Post Supply Chain,''
\emph{IEEE Access}, vol.~5, pp.~17465--17477, 2017.

\bibitem{azaria2016medrec}
A.~Azaria, A.~Ekblaw, T.~Vieira, and A.~Lippman,
``MedRec: Using Blockchain for Medical Records, Decentralized
Content Management,''
in \emph{Proc.\ IEEE Int.\ Conf.\ Open and Big Data}, 2016,
pp.~25--30.

\bibitem{shapley1954method}
L.~S.~Shapley and M.~Shubik,
``A Method for Evaluating the Distribution of Power in a Committee
System,''
\emph{Amer.\ Polit.\ Sci.\ Rev.}, vol.~48, no.~3, pp.~787--792,
1954.

\bibitem{ranger}
Apache Software Foundation,
``Apache Ranger,'' 2023. \url{https://ranger.apache.org}

\bibitem{opa}
Open Policy Agent,
``OPA: Policy-based control for cloud native environments,'' 2023.
\url{https://www.openpolicyagent.org}

\bibitem{w3cdid}
W3C, ``Decentralized Identifiers (DIDs) v1.0,''
W3C Recommendation, 2022.

\bibitem{w3cvc}
W3C, ``Verifiable Credentials Data Model v2.0,''
W3C Recommendation, 2024.

\bibitem{popa2011cryptdb}
R.~A.~Popa, C.~M.~S.~Redfield, N.~Zeldovich, and H.~Balakrishnan,
``CryptDB: Protecting Confidentiality with Encrypted Query
Processing,''
in \emph{Proc.\ SOSP}, 2011, pp.~85--100.

\bibitem{naveed2015inference}
M.~Naveed, S.~Kamara, and C.~V.~Wright,
``Inference Attacks on Property-Preserving Encrypted Databases,''
in \emph{Proc.\ CCS}, 2015, pp.~644--655.

\bibitem{rogaway2006siv}
P.~Rogaway and T.~Shrimpton,
``Deterministic Authenticated-Encryption: A Provable-Security
Treatment of the Key-Wrap Problem,'' in \emph{Proc.\ EUROCRYPT},
2006, pp.~373--390.

\bibitem{rizvi2004nontruman}
S.~Rizvi, A.~Mendelzon, S.~Sudarshan, and P.~Roy,
``Extending Query Rewriting Techniques for Fine-Grained Access
Control,'' in \emph{Proc.\ SIGMOD}, 2004, pp.~551--562.

\end{thebibliography}

\end{document}
